\runninghead{Three Interpretations Compared}
\section{Where the Three Interpretations Agree and Differ}
\label{sec:comparison}

The three interpretations take up the three questions set out at the opening,
and each starts from different texts of the Mass.

\begin{comparisontable}{}{One Church, sealed and uncounted}{The way of the Beatitudes}{No want to them that fear him}
Question &
Who are the sealed, counted at 144,000 and then shown as a multitude no one
could number? &
Why does the feast read the Beatitudes, which bless the poor and the
persecuted, and how are their present and future promises to be read? &
How can the Gradual promise that nothing is lacking to those who fear the
Lord, and the Offertory say the just only seemed to die? \\
Texts it starts from &
Lesson, Introit verse &
Gospel, Communion, Alleluia &
Gradual, Offertory \\
Answer &
The number signifies the whole Church, sealed with the cross; the multitude
is the victors of the last persecution or the Gentiles added to Israel &
The beatitudes are one way, from humility and fear to persecution for
justice, whose one reward is begun on earth and completed in heaven &
What the God-fearing are promised is God himself and the virtues; want, when
it comes, is sent for their health; the just dead are in God's hand \\
Principal authors &
Bede, with Augustine on the psalms &
Augustine, Chrysostom; Aquinas, Hilary &
Basil, Bellarmine; Augustine, Aquinas \\
What it asks of the faithful &
To confess the faith openly and rejoice in the Lord, not in themselves &
To begin in humility and the fear of God and be ready to suffer for justice &
Filial fear, and patience in want \\
Strongest difficulty &
The exact number and Israel's tribes, read as a limit &
The tenses: two promises present, six future &
The promise read as temporal comfort, which saints' lives contradict \\
Where its witnesses differ &
Bede reads the number as the whole Church; Fromage reads the sealing as a
second, definitive census; Schuster takes the 144,000 as sons of Abraham with
others added &
Augustine's ascent of knowledge and gifts against Chrysostom's chain of
dispositions; Hilary's land as the Lord's body against Chrysostom's sensible
prize &
Basil reads the rich as Israel, Bellarmine as the wicked rich and the
covetous; only Bellarmine keeps a temporal sense \\
\end{comparisontable}

All three rest on the fear of the Lord and on humility, though they come to
it from different texts. In the first, the just rejoice ``not in yourselves''
but in the Lord, and the seal is set on those who confess. In the second, the
poor in spirit are the humble and God-fearing, and Augustine's way begins
with the gift of fear. In the third, the Gradual's \latin{Timete Dominum} is
the subject itself, and Basil and Bellarmine read that fear as the good that
lacks nothing. The three agree that the saints are not honoured on this feast
for what they achieved by themselves.

They attend to different times. The first looks at the saints already
before the throne and reads back from there to the sealing on earth. The
second follows the way the saints walked, from its first step to its last.
The third stands where the assembly stands, with the promise still being
tested by hunger and death. Bede's remark that the vision of the white-robed
multitude ``may also be understood of the present time, when `we are saved in
hope'\,'' and
Aquinas's distinction between the rewards begun on the way and completed in
the homeland, give the first two a bridge to the third's present.

The first and the third read the Lesson and the Offertory differently. For
the first, the multitude with palms is the Church that has overcome, and the
Offertory's just are the sealed who have passed through death. For the third,
the Offertory answers the scandal of the just man's death, and the Lesson is
not its seat: it rests on the psalm and on Wisdom, not on the Apocalypse. The
second and the third share the Gospel's fourth beatitude, the hunger and
thirst after justice, which Augustine cites inside his exposition of the
Gradual's psalm; for the second it is a step on the way, for the third it
names the want that is blessed.

Each accounts for texts that the others pass over. Only the first explains
the Lesson's number and its list of tribes, and it does so with Bede's reasons
rather than by dropping the number. Only the second explains why the Gospel
is the Beatitudes and why their promises are in two tenses. Only the third
meets the objection that the Gradual's promise is false, and says what the
Mass sings and what it leaves out of Psalm~33. The questions differ, and the
answers do not contradict one another. Where their witnesses disagree, on the
census, on the land of the meek and on the rich who lack, the study has kept
both sides.
